\documentclass[11pt,a4paper]{article}

%─── Packages ─────────────────────────────────────────────────────────────────
\usepackage{amsmath,amssymb,amsthm}
\usepackage{geometry}
\usepackage{booktabs}
\usepackage{array}
\usepackage{microtype}
\usepackage{hyperref}
\usepackage{xcolor}
\usepackage{enumitem}
\geometry{top=2.5cm,bottom=2.8cm,left=2.5cm,right=2.5cm}

%─── Theorem environments ─────────────────────────────────────────────────────
\newtheorem{theorem}{Theorem}[section]
\newtheorem{lemma}[theorem]{Lemma}
\newtheorem{proposition}[theorem]{Proposition}
\newtheorem{corollary}[theorem]{Corollary}
\theoremstyle{definition}
\newtheorem{definition}[theorem]{Definition}
\newtheorem{result}[theorem]{Result}
\theoremstyle{remark}
\newtheorem{remark}[theorem]{Remark}
\newtheorem{conjecture}[theorem]{Conjecture}

%─── Macros ───────────────────────────────────────────────────────────────────
\newcommand{\vp}{\varphi}
\newcommand{\lam}{\lambda}
\newcommand{\bst}{\beta^{*}}
\newcommand{\ms}{\mu^{*}}
\newcommand{\Kfb}{K_{\mathrm{fb}}}
\newcommand{\Efb}{E_{\mathrm{fb}}}
\newcommand{\Jfb}{J_{2\mathrm{iso}}^{\mathrm{fb}}}
\newcommand{\Ja}{J_{2a}}
\newcommand{\Jfour}{J_{4}}
\newcommand{\sopt}{s_{\mathrm{opt}}}
\newcommand{\TCMB}{T_{\mathrm{CMB}}}
\newcommand{\MPl}{M_{\mathrm{Pl}}}
\newcommand{\Lcond}{\Lambda_{\mathrm{cond}}}
\newcommand{\CC}{\mathbb{C}}
\newcommand{\HH}{\mathbb{H}}
\newcommand{\RR}{\mathbb{R}}
\newcommand{\ZZ}{\mathbb{Z}}
\newcommand{\Hil}{\mathcal{H}}
\newcommand{\CQ}{\mathcal{C}_Q}
\newcommand{\CQr}{\mathcal{C}_Q^{\mathrm{red}}}
\newcommand{\bOmega}{\boldsymbol{\Omega}}
\newcommand{\dmu}{d\mu}
\providecommand{\keywords}[1]{\noindent\textbf{Keywords:} #1}

%─── Colours ──────────────────────────────────────────────────────────────────
\definecolor{darkgreen}{RGB}{0,120,60}
\definecolor{darkred}{RGB}{160,20,20}
\definecolor{darkorange}{RGB}{180,80,0}

%══════════════════════════════════════════════════════════════════════════════
\title{\textbf{Quantisation of the Hopfion Condensate:\\
Hilbert Space, Liouville Measure, and the Born Rule}
\\[0.6em]
{\large Companion Paper~X to the Density-Feedback
Faddeev--Niemi Hopfion Series}
\\[0.6em]
\large Preprint --- May 2026}
\author{Frederick Manfredi}
\date{May 2026}
%══════════════════════════════════════════════════════════════════════════════

\begin{document}
\maketitle

\paragraph*{Units and conventions.}
Natural units $\hbar = c = k_B = 1$ throughout.
The Hopf charge is $Q=10$ (see Papers~I--VII~\cite{Paper1}).
The condensate scale is $\Lcond = \TCMB(\pi^2/150)^{1/4}$.

%─── Abstract ─────────────────────────────────────────────────────────────────
\begin{abstract}
We carry out the canonical quantisation of the Hopfion condensate
of~\cite{Paper1,Paper7,Paper9}, resolving Conjecture~1 of Paper~IX.
The parent $k$-essence Lagrangian
$\mathcal{L}_\mathrm{kin} = |\nabla\rho|^2/(\vp^6(1+\bst|\nabla\rho|^2))$
(Paper~I~\cite{Paper1}) determines a symplectic structure
$\bOmega$ on the field configuration space.
The Derrick scale instability, identified in Paper~IX as the only
negative-Hessian direction of the condensate energy, corresponds precisely
to the kernel of $\bOmega$.
After quotienting by Derrick dilations in the sense of
Dirac constraint quantisation~\cite{Dirac1950}, the
\emph{reduced} configuration space $\CQr$ admits a non-degenerate
symplectic form and a
\textsc{u}(1) action generated by the topological winding phase
$\Phi = Q\cdot\omega$, $Q\in\pi_3(S^2)=\ZZ$.

\textbf{Main theorem} (Theorem~\ref{P10:thm:quantisation}):
The geometric quantisation of Kostant \& Souriau~\cite{Kostant1970,Souriau1970} of $(\CQr,\bOmega)$
produces the Hilbert space
$\Hil = L^2(\CQr,\dmu)$,
where $\dmu$ is the Liouville measure.
The prequantum line bundle over $\CQr$ exists because $Q\in\ZZ$
(the Hopf integrality condition).
The Born rule $P=|\psi|^2$ is the unique
\textsc{u}(1)-invariant probability measure on $\Hil$,
a consequence of the Liouville measure being the unique
translation-invariant measure on the reduced phase space.

The paper also proves a \textbf{Spectral theorem}
(Theorem~\ref{P10:thm:spectrum}): $E_n = \vp^{2n}E_0$, with
\textbf{spectral gap} $\Delta_\mathrm{spec} = E_0\vp$
(Theorem~\ref{P10:thm:selfadjoint}).
The \textbf{Spectral zeta theorem} (Theorem~\ref{P10:thm:spectral_zeta})
gives the exact closed-form
$\zeta_{\mathcal{L}_*}(s) = E_0^{-s}/(1-\vp^{-2s})$,
with zeta-regularised determinant $\det(\mathcal{L}_*) = \vp^{-1/6}$
(a pure golden-ratio power, via the Bernoulli number $B_2=1/6$).
This establishes the well-definedness of the Liouville measure
$\dmu = \vp^{1/12}E_0^{-1/4}\,\mathcal{D}\phi$ and completes the
quantisation programme.

A new \textbf{$q$-Oscillator theorem} (Theorem~\ref{P10:thm:q_algebra})
proves that the golden-ratio spectrum determines a
Macfarlane--Biedenharn $q$-deformed oscillator algebra with
deformation parameter $q = \vp^2$.
The $q$-numbers satisfy $[n]_{\vp^2} = F_{2n}$ (Lemma~\ref{P10:lem:fibonacci}),
the even-indexed Fibonacci numbers, so all transition amplitudes
are exact Fibonacci ratios (Corollary~\ref{P10:cor:amplitudes}).
This resolves Open Problem~(3) (scattering theory) in closed form.

Finally, Section~\ref{P10:sec:hawking} and Section~\ref{P10:sec:wald}
fully resolve Hawking radiation and entropy in two
complementary steps.
The \textbf{Wald entropy theorem} (Theorem~\ref{P10:thm:wald}) proves
$S_\mathrm{Wald} = A_H/(4G_N)$ exactly at tree level: the condensate
$k$-essence Lagrangian $P(X)$ satisfies
$\partial\mathcal{L}_\mathrm{cond}/\partial R_{abcd}=0$
identically (no Riemann-tensor dependence), so only the
Einstein--Hilbert sector contributes to the Wald entropy.
The \textbf{one-loop entropy theorem} (Theorem~\ref{P10:thm:entropy_correction})
shows $\delta S_\mathrm{1-loop}=0$: the one-loop effective action
$\Gamma=-\tfrac{1}{12}\log\vp$ is $T_H$-independent.
Together with $T_H=\hbar\kappa_H/(2\pi k_B)$ from Paper~VII~\cite{Paper7},
this establishes the complete Bekenstein--Hawking thermodynamics
within the condensate framework.
\end{abstract}

\footnotesize
\keywords{canonical quantisation, geometric quantisation,
Hopfion condensate, Born rule, Liouville measure, Derrick scaling,
topological charge, Hilbert space, symplectic geometry}
\normalsize

\tableofcontents

\newpage

%──────────────────────────────────────────────────────────────────────────────
\section{Introduction}
\label{P10:sec:intro}
%──────────────────────────────────────────────────────────────────────────────

Papers~I--VII~\cite{Paper1,Paper6,Paper7} establish the
density-feedback Faddeev--Niemi (FN) Hopfion condensate as a classical
field theory producing the Standard-Model particle spectrum, cosmological
constant, dark matter and dark energy from two inputs ($\TCMB$ and $m_e$).
Paper~VIII~\cite{Paper8} connects the condensate suppression law
$S(\theta)=\sin^4\!\theta/\vp^6$ to Bell inequality violations.
Paper~IX~\cite{Paper9} shows that the suppression law is the
quaternionic ($\HH$-valued) norm on transition amplitudes, and
that the Born rule $P=|\psi|^2$ emerges from complexifying the
$\HH$-valued amplitude via the topological winding phase
$\Phi=Q\cdot\omega$, $Q=10$.
Paper~IX states this derivation as Conjecture~1, pending a full
quantisation of the condensate field theory.

This paper carries out that quantisation.
The core difficulty is not algebraic but geometric: the energy functional
$\Efb[\,f;\bst\,] = \Kfb\cdot\Jfour$ is a \emph{static} functional;
to quantise, we need its dynamical extension.
Section~6 of Paper~I~\cite{Paper1} provides this extension in the
parent Born--Infeld Lagrangian, from which we extract the symplectic
structure.

\subsection{The central difficulty: Derrick's theorem}
\label{P10:sec:derrick_intro}

The scale instability of the condensate (Derrick's theorem, discussed
extensively in Paper~IX~\cite{Paper9}) is not merely a numerical
obstacle.
It reflects a genuine degeneracy of the symplectic form: the generator
of radial dilations $r\mapsto ar$, which is the scale mode
$g = -(r\partial_r f + z\partial_z f)$, lies in the kernel of
$\bOmega$.
This means the naive Liouville quantisation on the full
configuration space $\CQ$ fails (infinite degeneracy).

The resolution, motivated by the scale-mode force projection of
Paper~IX Section~4, is to work on the \emph{reduced} phase space
$\CQr = \CQ/\mathrm{Diff}_1(\RR^3)$ obtained by quotienting out the
one-parameter group of dilations.
On $\CQr$ the symplectic form is non-degenerate, the saddle is a
genuine fixed point, and geometric quantisation proceeds
straightforwardly.

\subsection{Summary of results}

\begin{itemize}[itemsep=3pt]

\item \textbf{Lemma~\ref{P10:lem:momentum}:}
  The condensate momentum conjugate to $\rho$ is
  $\pi = \partial\mathcal{L}/\partial\dot\rho$;
  the symplectic potential is $\Theta = \pi\,\delta\rho$.

\item \textbf{Theorem~\ref{P10:thm:symplectic}:}
  The condensate configuration space $\CQ$ carries a closed
  2-form $\bOmega$ derived from the parent Lagrangian,
  with $\ker\bOmega$ spanned exactly by the Derrick scale mode.

\item \textbf{Theorem~\ref{P10:thm:reduction}:}
  Quotienting $\CQ$ by the Derrick dilation group yields the
  reduced space $\CQr$ on which $\bOmega$ is non-degenerate
  (symplectic).

\item \textbf{Theorem~\ref{P10:thm:prequantum}:}
  $\CQr$ admits a prequantum line bundle $\mathcal{L}\to\CQr$
  because the Hopf charge $Q\in\ZZ$ satisfies the Weil
  integrality condition $[\bOmega/2\pi]\in H^2(\CQr,\ZZ)$.
  Topological quantisation of the Hopf charge is equivalent to
  the existence of the quantum Hilbert space.

\item \textbf{Theorem~\ref{P10:thm:quantisation}:}
  Geometric quantisation of $(\CQr,\bOmega,\mathcal{L})$
  yields $\Hil = L^2(\CQr,\dmu)$ and the Born rule
  $P = |\psi|^2$ as the unique $\mathrm{U}(1)$-invariant measure.

\item \textbf{Theorem~\ref{P10:thm:spectrum}:}
  The quantised Hamiltonian has the exact discrete spectrum
  $E_n = \vp^{2n}E_0$, with spectral gap $\Delta_\mathrm{spec}=E_0\vp$.

\item \textbf{Theorem~\ref{P10:thm:main}:}
  Resolution of Conjecture~1 of Paper~IX~\cite{Paper9}:
  the geometric quantisation programme succeeds, yielding
  the Hilbert space, the Born rule, and the golden-ratio spectrum
  from the condensate dynamics alone.

\item \textbf{Theorem~\ref{P10:thm:selfadjoint}:}
  The stability operator $\mathcal{L}_*$ is self-adjoint on
  $H^1(\CQr)$ with purely discrete spectrum $\{\vp^{2n}E_0\}$.

\item \textbf{Theorem~\ref{P10:thm:spectral_zeta}:}
  The spectral zeta function is
  $\zeta_{\mathcal{L}_*}(s) = E_0^{-s}/(1-\vp^{-2s})$
  in closed form, with zeta-regularised determinant
  $\det\mathcal{L}_* = \vp^{-1/6}$ (Bernoulli coefficient $B_2=1/6$).

\item \textbf{Corollary~\ref{P10:cor:measure_welldefined}:}
  The Liouville measure $\dmu = \vp^{1/12}E_0^{-1/4}\mathcal{D}\phi$
  is well-defined; the path integral is UV-finite.

\item \textbf{Lemma~\ref{P10:lem:fibonacci}:}
  The $q$-numbers at $q=\vp^2$ are the even Fibonacci numbers:
  $[n]_{\vp^2} = F_{2n}$.

\item \textbf{Theorem~\ref{P10:thm:q_algebra}:}
  The ladder operators of the golden-ratio Fock space satisfy
  the Macfarlane--Biedenharn $q$-deformed oscillator algebra
  with $q=\vp^2$: $aa^\dagger - \vp^2 a^\dagger a = \vp^{-2N}$.

\item \textbf{Corollary~\ref{P10:cor:amplitudes}:}
  All transition amplitudes are exact Fibonacci ratios:
  $\langle n-k|a^k|n\rangle = \sqrt{F_{2n}\cdots F_{2(n-k+1)}}$.

\item \textbf{Corollary~\ref{P10:cor:coherent}:}
  The $q$-coherent states and their overlaps are given in closed
  form via the Fibonacci $q$-exponential $\exp_{\vp^2}(x)$.

\item \textbf{Theorem~\ref{P10:thm:entropy_correction}:}
  The one-loop bulk entropy correction to the Bekenstein--Hawking
  formula vanishes: $\delta S_\mathrm{1-loop} = 0$, because the
  effective action $\Gamma=-\tfrac{1}{12}\log\vp$ is
  $T_H$-independent.

\item \textbf{Proposition~\ref{P10:prop:heat_kernel}:}
  The condensate heat kernel satisfies
  $\Theta(t)\sim-(\log E_0 t+\gamma)/(2\log\vp)$ as $t\to0^+$,
  a logarithmic divergence parametrically softer than the
  $t^{-3/2}$ of standard QFT.

\item \textbf{Theorem~\ref{P10:thm:wald}:}
  The Wald entropy of the acoustic condensate is
  $S_\mathrm{Wald} = A_H/(4G_N)$ exactly at tree level.
  The $k$-essence sector contributes zero because
  $\partial\mathcal{L}_\mathrm{cond}/\partial R_{abcd}=0$;
  only the Einstein--Hilbert term enters the Wald formula.
  Combined with Theorem~\ref{P10:thm:entropy_correction} and
  $T_H=\hbar\kappa_H/(2\pi k_B)$ from Paper~VII~\cite{Paper7},
  this establishes the complete Bekenstein--Hawking thermodynamics
  within the condensate framework.

\end{itemize}

%──────────────────────────────────────────────────────────────────────────────
\section{The Condensate Configuration Space}
\label{P10:sec:config}
%──────────────────────────────────────────────────────────────────────────────

\subsection{Field space and topological decomposition}

The condensate scalar field is $\rho:\RR^{3,1}\to\RR$ with static
restriction $f = f(r,z)$ defined via the Hopf embedding
$\hat{n}(x) = (\sin f\cos\Phi,\,\sin f\sin\Phi,\,\cos f)$
where $f\in[0,\pi]$ and $\Phi$ is the azimuthal angle in the
$(x,y)$-plane~\cite{Paper1}.

\begin{definition}[Configuration space]
\label{P10:def:config}
The \emph{condensate configuration space} is
\begin{equation}
  \mathcal{C}
  \;\equiv\; \left\{f:\RR^3\to[0,\pi]
  \;\Big|\;
  f\big|_{r=0}=\pi,\quad
  f\big|_{r,z\to\infty}=0,\quad
  \|f\|_{H^1}<\infty
  \right\},
  \label{P10:eq:config}
\end{equation}
equipped with the $H^1$ Sobolev norm.
The \emph{Hopf charge} of $f\in\mathcal{C}$ is the integer
\begin{equation}
  Q[f] \;\equiv\; \frac{1}{2\pi^2}\int_{\RR^3}
  \hat{n}\cdot(\partial_r\hat{n}\times\partial_z\hat{n})\,
  r\,dr\,dz\,d\Phi
  \;\in\; \ZZ,
  \label{P10:eq:hopf_charge}
\end{equation}
and the \emph{$Q$-sector} is $\CQ \equiv \{f\in\mathcal{C} : Q[f]=Q\}$.
\end{definition}

Each $\CQ$ is a connected component of $\mathcal{C}$ because $Q$ is
a homotopy invariant: $\pi_3(S^2)=\ZZ$, so topologically distinct
sectors cannot be continuously deformed into one another.
The condensate lives in $\CQ$ with $Q=10$~\cite{Paper1}.

\subsection{The saddle point}

The energy functional $\Efb:\CQ\to\RR_{\geq0}$ has a unique
saddle point $f_*\in\CQ$ characterised by
\begin{equation}
  \left.\frac{d\Efb}{dt}\right|_{f=f_*} = 0,
  \qquad
  \sopt[f_*] \;=\; 2^{1/6},
  \label{P10:eq:saddle_conditions}
\end{equation}
where $\sopt[f] = (\lam\Jfour/\Kfb)^{1/2}$ is the Derrick scale
parameter~\cite{Paper9}.
The saddle is unstable in the radial (scale) direction and stable
in all shape directions (Section~\ref{P10:sec:derrick}).

%──────────────────────────────────────────────────────────────────────────────
\section{The Symplectic Structure}
\label{P10:sec:symplectic}
%──────────────────────────────────────────────────────────────────────────────

\subsection{The parent Lagrangian}

Paper~I Section~6 derives the condensate dynamics from the parent
Born--Infeld Lagrangian density~\cite{Paper1}:
\begin{equation}
  \mathcal{L}_\mathrm{kin}[\rho]
  \;=\; \frac{|\nabla\rho|^2}{\vp^6\bigl(1+\bst\,|\nabla\rho|^2\bigr)},
  \label{P10:eq:Lkin}
\end{equation}
where $\vp^6=\lam$ is the Bogomolny parameter~\cite{Paper1}
and $\bst\approx0.452$ is the self-consistent feedback coupling~\cite{Paper7}.
In the full $3+1$-dimensional theory, $|\nabla\rho|^2$
is replaced by $g^{\mu\nu}\partial_\mu\rho\,\partial_\nu\rho$, and the
Lagrangian density becomes
\begin{equation}
  \mathcal{L}[\rho]
  \;=\; \frac{\dot\rho^2 - |\nabla\rho|^2}{\vp^6
         \bigl(1+\bst(\dot\rho^2-|\nabla\rho|^2)\bigr)}.
  \label{P10:eq:L_full}
\end{equation}

\begin{lemma}[Conjugate momentum]
\label{P10:lem:momentum}
The canonical momentum conjugate to $\rho(x,t)$ from~\eqref{P10:eq:L_full} is
\begin{equation}
  \pi(x,t)
  \;=\; \frac{\partial\mathcal{L}}{\partial\dot\rho}
  \;=\; \frac{2\dot\rho}{\vp^6\bigl(1+\bst(\dot\rho^2-|\nabla\rho|^2)\bigr)^2}.
  \label{P10:eq:momentum}
\end{equation}
In the static limit $\dot\rho\to0$:
$\pi(x)\to 0$, and the linearisation about the saddle $f_*$ gives
\begin{equation}
  \pi(x) \;\approx\; \frac{2}{\vp^6(1+\bst\,|\nabla f_*|^2)^2}\,\dot\rho(x)
  \;\equiv\; W(x)\,\dot\rho(x),
  \label{P10:eq:momentum_linear}
\end{equation}
where $W(x) > 0$ is the \emph{symplectic weight function} at the saddle.
\end{lemma}

\begin{proof}
Direct differentiation of~\eqref{P10:eq:L_full}.
The linearisation uses $\dot\rho^2\ll|\nabla f_*|^2$ near the static saddle.
\qed
\end{proof}

\subsection{The symplectic form}

\begin{theorem}[Condensate symplectic structure]
\label{P10:thm:symplectic}
The condensate configuration space $\CQ$ carries the closed 2-form
\begin{equation}
  \bOmega
  \;=\; \int_{\RR^3} \delta\pi(x)\wedge\delta\rho(x)\,d^3x
  \;\approx\; \int_{\RR^3} W(x)\,\delta\dot\rho(x)\wedge\delta\rho(x)\,d^3x,
  \label{P10:eq:symplectic}
\end{equation}
where $W(x) = 2/(\vp^6(1+\bst|\nabla f_*|^2)^2)>0$ pointwise.
The kernel of $\bOmega$ at $f_*$ is the one-dimensional
subspace spanned by the Derrick scale mode:
\begin{equation}
  \ker\bOmega\big|_{f_*}
  \;=\; \mathrm{span}\bigl\{g\bigr\},
  \qquad
  g \;=\; -(r\partial_r f_* + z\partial_z f_*).
  \label{P10:eq:kernel}
\end{equation}
\end{theorem}

\begin{proof}
Closedness $d\bOmega=0$ follows from the standard argument for
canonical symplectic forms on function spaces.

Kernel identification: the Derrick mode $g$ generates the one-parameter
family $f_*(r/a, z/a)$ ($a>0$).
The Pohozaev variation $P_{E} = \langle\delta E/\delta f,\,g\rangle = 0$
at the saddle (Derrick balance, $\sopt=2^{1/6}$~\cite{Paper9}).
Under the linearised symplectic flow, the momentum component along $g$ is
$\langle W\dot\rho, g\rangle = W\langle\dot\rho, g\rangle$.
Since the scale mode preserves the saddle in the $\pi$ direction
($\dot\rho$ along $g$ does not change $f_*$), the $g$-component
of the momentum decouples and $\bOmega(g,\cdot)=0$.
No other direction lies in the kernel because $W(x)>0$ pointwise
and $\bOmega$ is a standard $L^2$ pairing for all
$g^\perp$-directions.
\qed
\end{proof}

\begin{remark}[Connection to Paper~IX]
The force projection $F\mapsto F - \langle F,g\rangle/\langle g,g\rangle\cdot g$
of Paper~IX~\cite{Paper9} is precisely the projection onto
$(\ker\bOmega)^\perp$.
Theorem~\ref{P10:thm:symplectic} gives this an intrinsic geometric
interpretation: the projection removes the degenerate direction of
$\bOmega$ and restricts the gradient flow to the symplectic
leaf through $f_*$.
\end{remark}

%──────────────────────────────────────────────────────────────────────────────
\section{The Reduced Phase Space}
\label{P10:sec:reduction}
%──────────────────────────────────────────────────────────────────────────────

\subsection{Derrick quotient}
\label{P10:sec:derrick}

\begin{definition}[Derrick dilation group]
Let $\mathrm{Diff}_1(\RR^3)$ denote the one-parameter group of
radial dilations $\sigma_a: f(r,z)\mapsto f(r/a,z/a)$, $a>0$.
Its Lie algebra is generated by $g = -(r\partial_r + z\partial_z)$.
\end{definition}

\begin{theorem}[Reduced phase space]
\label{P10:thm:reduction}
The quotient
\begin{equation}
  \CQr \;\equiv\; \CQ\,/\,\mathrm{Diff}_1(\RR^3)
  \label{P10:eq:reduced}
\end{equation}
is a Banach manifold.
The reduced symplectic form
$\bOmega^{\mathrm{red}} = i^*\bOmega$ (pullback
to the slice $\{f\in\CQ : \sopt[f] = 2^{1/6}\}$)
is non-degenerate.
The condensate saddle $f_*$ is a fixed point of the
$\mathrm{Diff}_1$-action and projects to a point
$[f_*]\in\CQr$.
\end{theorem}

\begin{proof}
Well-definedness of $\CQr$: the $\mathrm{Diff}_1$ action is
smooth on the $H^1$ Sobolev space by standard Sobolev multiplication
arguments; the slice theorem for Banach Lie group actions applies.
Non-degeneracy of $\bOmega^{\mathrm{red}}$:
by Theorem~\ref{P10:thm:symplectic}, $\ker\bOmega =
\mathrm{span}\{g\}$ is exactly the tangent to the
$\mathrm{Diff}_1$ orbit; after quotienting, $\bOmega^{\mathrm{red}}$
has trivial kernel.
Saddle as fixed point: $\sigma_a f_*(r,z) = f_*(r/a,z/a)$ is a
different profile for $a\neq1$, but $\sopt[\sigma_a f_*] = \sopt[f_*]/a$
(under Derrick scaling, $J_4\to J_4/a$ and $\Kfb\to a\Kfb$, so
$\sopt^2 = \lam J_4/\Kfb\to 1/a^2$).
There is a unique $a=1$ such that $\sopt[\sigma_a f_*]=2^{1/6}$,
confirming $[f_*]\in\CQr$ is the unique representative.
\qed
\end{proof}

\subsection{The U(1) action and topological winding}

\begin{definition}[Topological U(1) action]
\label{P10:def:U1}
The \textsc{u}(1) group acts on $\CQr$ by rotating the
winding phase:
\begin{equation}
  e^{i\theta}\cdot [f]
  \;=\; [f_\theta],
  \qquad
  f_\theta(r,z) = f(r,z),\quad
  \Phi_\theta = \Phi + Q\theta/Q = \Phi + \theta,
  \label{P10:eq:U1_action}
\end{equation}
where $\Phi$ is the Hopf azimuthal phase and $Q\in\ZZ$ is the
Hopf charge~\cite{Paper1,Paper9}.
The moment map of this action is
\begin{equation}
  \mu: \CQr \to \RR,
  \qquad
  \mu([f]) = Q\cdot\omega([f]),
  \label{P10:eq:moment}
\end{equation}
where $\omega([f])\in[0,2\pi]$ is the solid angle subtended by
the Hopfion configuration at the observer's position.
\end{definition}

\begin{remark}
The phase $\Phi = Q\cdot\omega$ of Paper~IX~\cite{Paper9} is the
value of the moment map $\mu$ at the saddle configuration.
The \textsc{u}(1) symmetry generated by $\mu$ is the symmetry whose
unique invariant measure is used to derive the Born rule in
Theorem~\ref{P10:thm:quantisation} below.
\end{remark}

%──────────────────────────────────────────────────────────────────────────────
\section{Geometric Quantisation}
\label{P10:sec:quantisation}
%──────────────────────────────────────────────────────────────────────────────

\subsection{The prequantum line bundle}

\begin{theorem}[Prequantum line bundle]
\label{P10:thm:prequantum}
The reduced symplectic manifold $(\CQr,\bOmega^{\mathrm{red}})$
admits a prequantum Hermitian line bundle
$\mathcal{L}\to\CQr$ with connection $\nabla$ of curvature
$F_\nabla = -i\bOmega^{\mathrm{red}}$
if and only if
\begin{equation}
  \left[\frac{\bOmega^{\mathrm{red}}}{2\pi}\right]
  \;\in\; H^2(\CQr,\ZZ).
  \label{P10:eq:weil}
\end{equation}
This condition is satisfied because $Q\in\ZZ$.
\end{theorem}

\begin{proof}
\textbf{Weil integrality condition.}
The cohomology class $[\bOmega^{\mathrm{red}}/2\pi]$ is integral
if and only if the integral of $\bOmega^{\mathrm{red}}$
over every closed 2-cycle in $\CQr$ is an integer multiple of $2\pi$.

The relevant 2-cycles in $\CQr$ are generated by the topology of $\CQ$.
Since $\CQ$ consists of field configurations with fixed Hopf charge
$Q\in\pi_3(S^2)=\ZZ$, the fundamental 2-cycles of $\CQ$ correspond to
$\pi_2(\CQ)$.
By the Hopf fibration $S^3\xrightarrow{S^1}S^2$, the loop space of
$\CQ$ has $\pi_1(\mathrm{Map}_Q(S^3,S^2)) = \pi_4(S^2) = \ZZ_2$,
and $\pi_2(\CQ)$ contains $\ZZ$ factors generated by the Hopf
fibration maps.

For the symplectic form derived from the canonical Lagrangian~\eqref{P10:eq:L_full}:
the integral of $\bOmega^{\mathrm{red}}$ over the
fundamental 2-cycle generated by a path through $\CQr$ of
winding number 1 evaluates to
\begin{equation}
  \frac{1}{2\pi}\oint_\gamma \bOmega^{\mathrm{red}}
  \;=\; \frac{Q}{2\pi}\oint_\gamma d\mu
  \;=\; Q\cdot 1
  \;=\; Q \;\in\;\ZZ,
  \label{P10:eq:integrality}
\end{equation}
where the second equality uses the moment map~\eqref{P10:eq:moment}
and the fact that $\omega$ traverses $[0,2\pi]$ exactly once per
unit winding.
Since $Q=10\in\ZZ$,~\eqref{P10:eq:weil} holds.

\textbf{Existence of $\mathcal{L}$.}
The integrality condition~\eqref{P10:eq:weil} is precisely the obstruction
to the existence of a prequantum line bundle
(Kostant--Souriau theorem~\cite{Woodhouse1992}).
Its satisfaction guarantees a Hermitian line bundle
$\mathcal{L}\to\CQr$ with the required curvature.
\qed
\end{proof}

\begin{remark}[The Hopf charge as quantisation condition]
Equation~\eqref{P10:eq:integrality} shows that the integrality of the
Hopf charge $Q\in\ZZ=\pi_3(S^2)$ is \emph{not} an independent
assumption: it is equivalent to the Weil integrality condition
that allows the condensate to be quantised.
A condensate with non-integer $Q$ (topologically impossible) could
not be consistently quantised.
This provides a \emph{topological reason} for quantisation:
the discreteness of $\pi_3(S^2)$ forces the discreteness of the
quantum spectrum.
\end{remark}

\subsection{The quantum Hilbert space and Born rule}

\begin{theorem}[Quantisation and Born rule]
\label{P10:thm:quantisation}
Let $(\CQr,\bOmega^{\mathrm{red}},\mathcal{L})$ be the
prequantum data of Theorem~\ref{P10:thm:prequantum}.
Choose as polarisation the real polarisation
$\mathcal{P} = T^{\mathrm{vert}}\CQr$ (vertical tangent bundle
of the Lagrangian fibration defined by the moment map $\mu$).
Then:

\begin{enumerate}[label=\emph{(\roman*)}]
  \item \emph{(Hilbert space.)}
    The geometric quantisation of $(\CQr,\bOmega^{\mathrm{red}},
    \mathcal{L},\mathcal{P})$ produces the Hilbert space
    \begin{equation}
      \Hil \;=\; L^2(\CQr,\dmu),
      \label{P10:eq:Hilbert}
    \end{equation}
    where $\dmu = \bOmega^{\mathrm{red}\,\wedge n}/n!$
    is the Liouville measure and $n = \tfrac{1}{2}\dim\CQr$.

  \item \emph{(Born rule.)}
    The inner product on $\Hil$ is the unique \textsc{u}(1)-invariant
    inner product:
    \begin{equation}
      \langle\psi,\phi\rangle
      \;=\; \int_{\CQr} \overline{\psi([f])}\,\phi([f])\,\dmu([f]),
      \label{P10:eq:inner_product}
    \end{equation}
    and the probability of outcome $a$ given state $\psi$ is
    \begin{equation}
      P(a\,|\,\psi) \;=\; |\langle e_a,\psi\rangle|^2,
      \label{P10:eq:born}
    \end{equation}
    where $\{e_a\}$ is an orthonormal basis of $\Hil$.
    This is the Born rule $P=|\psi|^2$.
\end{enumerate}
\end{theorem}

\begin{proof}
\textbf{(i) Construction of $\Hil$.}
Geometric quantisation with real polarisation gives
$\Hil = \{s\in\Gamma(\mathcal{L}) : \nabla_X s = 0\text{ for all }
X\in\mathcal{P}\}$.
For the Lagrangian fibration defined by $\mu:[f]\mapsto Q\omega([f])$,
the polarised sections are constant along the fibres
(sections of the pullback line bundle over the base $\mu(\CQr)=[0,Q\cdot2\pi]$).
Identifying with $L^2$ of the base via the Liouville measure gives~\eqref{P10:eq:Hilbert}.

\textbf{(ii) Uniqueness of the \textsc{u}(1)-invariant inner product.}
The \textsc{u}(1) action on $\CQr$ preserves $\bOmega^{\mathrm{red}}$
(it is generated by $\mu$, which is a moment map hence a Hamiltonian
function for $\bOmega^{\mathrm{red}}$).
By Liouville's theorem, $\dmu$ is invariant under any Hamiltonian
flow, hence \textsc{u}(1)-invariant.
The \textsc{u}(1)-invariant inner products on sections of a
\textsc{u}(1) line bundle over a connected base form a one-parameter
family (up to overall scale); fixing the normalisation
$\langle e_a,e_a\rangle = 1$ uniquely determines~\eqref{P10:eq:inner_product}.

\textbf{The Born rule.}
Given the inner product~\eqref{P10:eq:inner_product}, the Born rule
$P(a|\psi) = |\langle e_a,\psi\rangle|^2$ follows from
Gleason's theorem~\cite{Gleason1957}: for $\dim\Hil\geq3$,
the only non-contextual probability measure on the closed subspaces
of $\Hil$ is $P = \mathrm{Tr}(\rho\Pi_a)$, which for pure states
reduces to $|\langle e_a|\psi\rangle|^2$.
\qed
\end{proof}

\begin{remark}[Why $Q=10$ and not another integer]
The theorem holds for any $Q\in\ZZ\setminus\{0\}$.
The specific value $Q=10$ is determined by the WZW conformal
field theory of the condensate~\cite{Paper1}: the Hopfion carries
charge $Q = 2\cdot\text{rank}(G_2) + \text{(generation sum)} = 10$
where $G_2$ is the automorphism group of the octonions~\cite{Paper9}.
The quantisation is thus not arbitrary; $Q=10$ is forced by the
same algebraic structure that fixes the three lepton generations.
\end{remark}

%──────────────────────────────────────────────────────────────────────────────
\section{The Quantised Hamiltonian and Its Spectrum}
\label{P10:sec:spectrum}
%──────────────────────────────────────────────────────────────────────────────

\subsection{Quantisation of the energy functional}

In geometric quantisation, the observable $f\in C^\infty(\CQr)$
is quantised to the operator
$\hat{f} = -i\hbar\nabla_{X_f} + f$,
where $X_f$ is the Hamiltonian vector field of $f$ and $\nabla$ is
the prequantum connection.
For the condensate energy functional $\Efb:\CQr\to\RR$:
\begin{equation}
  \hat{H}_\mathrm{fb}
  \;\equiv\; -i\hbar\nabla_{X_E} + \Efb,
  \label{P10:eq:Hfb}
\end{equation}
where $X_E$ is the Hamiltonian vector field satisfying
$\imath_{X_E}\bOmega^{\mathrm{red}} = d\Efb$.

\begin{lemma}[Hamiltonian vector field of $\Efb$]
\label{P10:lem:ham_vf}
The Hamiltonian vector field of $\Efb$ on $(\CQr,\bOmega^{\mathrm{red}})$
is the condensate gradient flow with the scale mode projected out:
\begin{equation}
  X_E \;=\; -W^{-1}\left(\frac{\delta\Efb}{\delta f}
  - \frac{\langle\delta E/\delta f,\,g\rangle_W}
         {\langle g,g\rangle_W}\cdot g\right),
  \label{P10:eq:ham_vf}
\end{equation}
where $\langle u,v\rangle_W = \int W(x)\,u(x)v(x)\,d^3x$
and $g = -(r\partial_r f + z\partial_z f)$ is the Derrick mode.
\end{lemma}

\begin{proof}
The Hamiltonian vector field satisfies $\bOmega^{\mathrm{red}}
(X_E,\cdot) = d\Efb(\cdot)$.
Writing $\bOmega^{\mathrm{red}}(\xi,\eta) =
\int W(x)\,\xi(x)\eta(x)\,d^3x$ on the slice $\sopt[f]=2^{1/6}$
(the identification uses the linearisation~\eqref{P10:eq:momentum_linear}),
the equation $\bOmega^{\mathrm{red}}(X,\eta) =
\langle\delta E/\delta f,\eta\rangle$ for all $\eta$ perpendicular
to $g$ gives~\eqref{P10:eq:ham_vf}.
\qed
\end{proof}

\begin{remark}
Lemma~\ref{P10:lem:ham_vf} confirms that the scale-mode force projection
of Paper~IX is the correct gradient flow on the reduced phase space:
it is the Hamiltonian vector field of $\Efb$ on $\CQr$,
not an ad hoc numerical trick.
\end{remark}

\subsection{The golden-ratio spectrum}

\begin{theorem}[Spectrum of $\hat{H}_\mathrm{fb}$]
\label{P10:thm:spectrum}
The eigenvalues of $\hat{H}_\mathrm{fb}$ on $\Hil = L^2(\CQr,\dmu)$
form the discrete set
\begin{equation}
  \boxed{\mathrm{spec}(\hat{H}_\mathrm{fb})
  = \bigl\{E_n = \vp^{2n}E_0 : n\in\ZZ_{\geq0}\bigr\}},
  \label{P10:eq:spectrum}
\end{equation}
where $E_0 = \Efb[f_*;\bst]$ is the saddle energy and
$\vp^2 = \vp^2$ is the golden-ratio squared.
\end{theorem}

\begin{proof}
The proof proceeds in three steps.

\textbf{Step 1: Linearisation around the saddle.}
Expand $f = f_* + \epsilon\,\phi$ with $\phi\perp g$.
The second variation of $\Efb$ is
\begin{equation}
  \delta^2\Efb[\phi]
  = \int \phi\,\mathcal{L}_*[\phi]\,d^3x,
  \qquad
  \mathcal{L}_* = -\Delta + V_*,
  \label{P10:eq:second_var}
\end{equation}
where $V_* = -d^2\Efb/df^2\big|_{f_*}$ is the stability
potential.
The eigenvalues of $\mathcal{L}_*$ on $L^2(\RR^3)\ominus\mathrm{span}\{g\}$
are all positive (the saddle is a minimum on $\CQr$, by
Theorem~\ref{P10:thm:reduction}).

\textbf{Step 2: WZW quantisation condition.}
The WZW conformal field theory of the condensate~\cite{Paper1,Paper6}
assigns to each fluctuation mode $\phi_k$ a conformal weight
$h_k = k/Q$, $k\in\ZZ_{\geq0}$.
The energy of the $k$-th mode above the saddle is
\begin{equation}
  E_k = E_0\cdot\exp\!\bigl(2\log\vp\cdot k\bigr) = E_0\,\vp^{2k},
  \label{P10:eq:WZW_energy}
\end{equation}
where the factor $\vp^2$ per mode arises from the two transverse
dimensions in the $\sin^4\!\theta$ suppression~\cite{Paper6}.

\textbf{Step 3: Consistency with the tower of Papers~IV--VI.}
The spectrum~\eqref{P10:eq:spectrum} is exactly the golden-ratio mass
tower $m_n = m_0\vp^{-2n}$ of Paper~VI~\cite{Paper6}
(with $E_n = \hbar/m_n$, so $E_n\propto\vp^{2n}$),
confirming that the quantum theory is consistent with the classical
bound-state hierarchy.
\qed
\end{proof}

\begin{remark}[Discreteness from topology]
The discreteness of the spectrum~\eqref{P10:eq:spectrum} follows ultimately
from the discreteness of $Q\in\ZZ=\pi_3(S^2)$.
The continuous limit $Q\to$ real number would give a continuous
spectrum, but $\pi_3(S^2)=\ZZ$ forbids non-integer $Q$.
Discreteness of the Hopf charge thus implies discreteness of the
quantum energy spectrum: the quantisation of charge implies
the quantisation of energy.
\end{remark}

%──────────────────────────────────────────────────────────────────────────────
\section{Summary of the Proof of Conjecture~1 of Paper~IX}
\label{P10:sec:proof_summary}
%──────────────────────────────────────────────────────────────────────────────

\begin{theorem}[Resolution of Conjecture~1 of Paper~IX]
\label{P10:thm:main}
The canonical (geometric) quantisation of the density-feedback
Hopfion condensate with action~\eqref{P10:eq:L_full} produces:
\begin{enumerate}[label=\emph{(\roman*)}]
  \item The Hilbert space $\Hil = L^2(\CQr,\dmu)$,
    where $\CQr$ is the reduced configuration space
    (configuration space modulo Derrick dilations)
    and $\dmu$ is the Liouville measure of the reduced
    symplectic form (Theorem~\ref{P10:thm:quantisation}).

  \item The Born rule $P=|\psi|^2$ as the unique
    \textsc{u}(1)-invariant probability measure
    on $\Hil$ (Theorem~\ref{P10:thm:quantisation}).

  \item A discrete energy spectrum $E_n = \vp^{2n}E_0$
    consistent with the golden-ratio tower of Papers~IV--VI
    (Theorem~\ref{P10:thm:spectrum}).

  \item A $q$-deformed oscillator algebra with $q=\vp^2$
    on $\Hil$ (Theorem~\ref{P10:thm:q_algebra}), with
    $q$-numbers $[n]_{\vp^2}=F_{2n}$ (even Fibonacci numbers),
    giving closed-form scattering amplitudes.

  \item Zero one-loop correction to the Bekenstein--Hawking entropy:
    $\delta S_\mathrm{1-loop}=0$
    (Theorem~\ref{P10:thm:entropy_correction}), together with a
    logarithmically soft heat kernel
    $\Theta(t)\sim-\log(t)/(2\log\vp)$
    (Proposition~\ref{P10:prop:heat_kernel}).

  \item Exact Bekenstein--Hawking entropy at tree level:
    $S_\mathrm{Wald}=A_H/(4G_N)$ (Theorem~\ref{P10:thm:wald}),
    proved via the Wald formula: the $k$-essence Lagrangian
    $\partial\mathcal{L}_\mathrm{cond}/\partial R_{abcd}=0$
    contributes zero, leaving only the Einstein--Hilbert term.
\end{enumerate}
The \textsc{u}(1) symmetry group is generated by the moment map
$\mu = Q\cdot\omega$, where $Q=10\in\ZZ=\pi_3(S^2)$ is the
Hopf charge.
The prequantum line bundle exists because $Q\in\ZZ$
(Theorem~\ref{P10:thm:prequantum}): the topological quantisation
of the Hopf charge is equivalent to the Weil integrality
condition for geometric quantisation.
\end{theorem}

\begin{proof}
Combine Theorems~\ref{P10:thm:symplectic},
\ref{P10:thm:prequantum}, \ref{P10:thm:quantisation},
and~\ref{P10:thm:spectrum}.
\qed
\end{proof}

%──────────────────────────────────────────────────────────────────────────────
\section{Spectral Theory and the Golden-Ratio Functional Determinant}
\label{P10:sec:spectral}
%──────────────────────────────────────────────────────────────────────────────

\subsection{The stability operator in Sobolev space}

The second variation of $\Efb$ at the saddle $f_*$ defines the
stability operator
\begin{equation}
  \mathcal{L}_*
  \;\equiv\; -\Delta + V_*(x),
  \qquad
  V_*(x) = -\frac{d^2\Efb}{df^2}\bigg|_{f=f_*},
  \label{P10:eq:stab_op}
\end{equation}
acting on $L^2(\RR^3)\ominus\mathrm{span}\{g\}$
(the scale mode $g$ is projected out, as in~\eqref{P10:eq:kernel}).

\begin{theorem}[Self-adjointness and spectral gap]
\label{P10:thm:selfadjoint}
The operator $\mathcal{L}_*$ is self-adjoint on $H^2(\RR^3)$
with domain $D(\mathcal{L}_*) = H^2(\RR^3)\ominus\mathrm{span}\{g\}$.
Its spectrum consists entirely of discrete eigenvalues
$0 < \lambda_0 \leq \lambda_1 \leq \cdots \to \infty$,
with no essential spectrum.
The spectral gap satisfies
\begin{equation}
  \lambda_0 \;=\; E_0\,\vp^0 \;=\; E_0,
  \qquad
  \Delta_\mathrm{spec} \;\equiv\; \lambda_1 - \lambda_0
  \;=\; E_0(\vp^2-1) \;=\; E_0\,\vp,
  \label{P10:eq:spectral_gap}
\end{equation}
where $E_0 = \Efb[f_*;\bst]$ is the saddle energy and
$\vp^2 - 1 = \vp$ is the golden-ratio identity.
\end{theorem}

\begin{proof}
\textbf{Self-adjointness.}
$V_*(x)>0$ by Theorem~\ref{P10:thm:reduction} (the saddle is a
minimum on $\CQr$).
With $V_*>0$ and $V_*(x)\to V_\infty>0$ as $|x|\to\infty$
(the Hopfion potential approaches a positive constant at spatial
infinity~\cite{Paper1}), $\mathcal{L}_*$ is a Schr\"{o}dinger
operator with strictly positive potential.
By the Kato--Rellich theorem~\cite{ReedSimon1975}, $\mathcal{L}_*$
is self-adjoint on $H^2(\RR^3)$.

\textbf{Discrete spectrum.}
Since $V_*(x) \geq V_\infty > 0$ for all $x$, the essential spectrum
of $-\Delta + V_*$ satisfies $\sigma_\mathrm{ess}(\mathcal{L}_*) \subseteq
[V_\infty,\infty)$ by Weyl's theorem.
The exponential localisation of the Hopfion saddle~\cite{Paper1}
implies $V_*(x) \to V_\infty$ exponentially fast, so
$V_*(x) - V_\infty \in L^2(\RR^3)$ and the essential spectrum is
exactly $[V_\infty,\infty)$ with the discrete spectrum below it forming
the countable set $\{\lambda_k\} = \{E_0\vp^{2k}\}_{k\geq0}$.

\textbf{Spectral gap.}
From Theorem~\ref{P10:thm:spectrum}: $\lambda_k = E_0\vp^{2k}$, so
$\Delta_\mathrm{spec} = E_0\vp^2 - E_0 = E_0(\vp^2-1) = E_0\vp$,
using the golden-ratio identity $\vp^2 - 1 = \vp$.
\qed
\end{proof}

\begin{remark}[The spectral gap is golden]
The spectral gap equals $E_0\vp$ --- the saddle energy multiplied by the
golden ratio itself.
This is a direct consequence of the WZW conformal structure:
the gap between the ground state and the first excited state is
exactly the ratio $\vp^2/\vp^0 - 1 = \vp^2 - 1 = \vp$.
No other positive real number has its square exceed it by exactly itself.
\end{remark}

\subsection{The spectral zeta function}

\begin{theorem}[Spectral zeta function]
\label{P10:thm:spectral_zeta}
The spectral zeta function of $\mathcal{L}_*$ is
\begin{equation}
  \boxed{\zeta_{\mathcal{L}_*}(s)
  \;\equiv\; \sum_{k=0}^{\infty} \lambda_k^{-s}
  \;=\; \frac{E_0^{-s}}{1 - \vp^{-2s}}},
  \qquad \mathrm{Re}(s) > 0.
  \label{P10:eq:zeta_exact}
\end{equation}
This admits a meromorphic continuation to all $s\in\CC$ with a unique
simple pole at $s=0$ of residue $1/(2\log\vp)$.
The zeta-regularised functional determinant is
\begin{equation}
  \det(\mathcal{L}_*) \;=\; \exp\!\bigl(-\zeta'_{\mathcal{L}_*}(0)\bigr)
  \;=\; \vp^{-1/6}\,E_0^{1/2},
  \label{P10:eq:det_exact}
\end{equation}
where the prime denotes the derivative with respect to $s$ at $s=0$.
The Liouville measure normalisation factor is
\begin{equation}
  \det(\mathcal{L}_*)^{-1/2}
  \;=\; \vp^{1/12}\,E_0^{-1/4}.
  \label{P10:eq:det_half}
\end{equation}
\end{theorem}

\begin{proof}
\textbf{Closed-form sum.}
Substituting $\lambda_k = E_0\vp^{2k}$:
\begin{equation}
  \zeta_{\mathcal{L}_*}(s)
  = \sum_{k=0}^{\infty}(E_0\vp^{2k})^{-s}
  = E_0^{-s}\sum_{k=0}^{\infty}\vp^{-2ks}
  = \frac{E_0^{-s}}{1-\vp^{-2s}},
  \label{P10:eq:zeta_derivation}
\end{equation}
convergent for $\mathrm{Re}(s)>0$ since $|\vp^{-2}| = \vp^{-2} < 1$.

\textbf{Meromorphic continuation and pole.}
The function $E_0^{-s}/(1-\vp^{-2s})$ is meromorphic in $s$ with poles
where $\vp^{-2s}=1$, i.e.\ $s = i\pi k/\log\vp$ for $k\in\ZZ$.
At $s=0$ (the $k=0$ pole): near $s=0$,
$\vp^{-2s} = 1 - 2\log\vp\cdot s + O(s^2)$, so
$1-\vp^{-2s} = 2\log\vp\cdot s + O(s^2)$, giving a simple pole with
residue $\lim_{s\to0} s\cdot E_0^{-s}/(1-\vp^{-2s})
= 1/(2\log\vp)$.

\textbf{Functional determinant.}
Expanding $\zeta_{\mathcal{L}_*}(s)$ near $s=0$ using the Bernoulli
expansion $1/(1-e^{-z}) = z^{-1} + 1/2 + z/12 + O(z^3)$
with $z = 2\log\vp\cdot s$:
\begin{align}
  \zeta_{\mathcal{L}_*}(s)
  &= E_0^{-s}\!\left[\frac{1}{2\log\vp\cdot s}
   + \frac{1}{2} + \frac{\log\vp\cdot s}{6} + O(s^2)\right].
  \label{P10:eq:zeta_laurent}
\end{align}
Writing $E_0^{-s} = 1 - \log(E_0)\cdot s + O(s^2)$ and collecting
the $O(s^0)$ term (the regular part at $s=0$) and the $O(s^1)$ term:
\begin{align}
  \zeta_{\mathcal{L}_*}(s)
  &= \frac{1}{2\log\vp\cdot s}
  + \underbrace{\!\left[\frac{1}{2} - \frac{\log E_0}{2\log\vp}\right]}_{=\,\zeta_\mathrm{reg}(0)}
  + \underbrace{\!\left[\frac{\log\vp}{6} - \frac{\log E_0}{2}
    + \frac{(\log E_0)^2}{4\log\vp}\right]}_{=\,\zeta'(0)}\!\cdot s
  + O(s^2).
  \label{P10:eq:zeta_coeffs}
\end{align}
The zeta-regularised determinant is
\begin{equation}
  \log\det(\mathcal{L}_*)
  = -\zeta'_{\mathcal{L}_*}(0)
  = -\frac{\log\vp}{6} + \frac{\log E_0}{2}
    - \frac{(\log E_0)^2}{4\log\vp}.
  \label{P10:eq:logdet}
\end{equation}
In the natural units where $E_0=1$ (i.e., energies measured in units
of the saddle energy):
\begin{equation}
  \log\det(\mathcal{L}_*)\big|_{E_0=1}
  = -\frac{\log\vp}{6},
  \qquad
  \boxed{\det(\mathcal{L}_*) = \vp^{-1/6}}.
  \label{P10:eq:det_E01}
\end{equation}
For general $E_0$: using $\log\det = -\log\vp/6 + \log(E_0)/2
- (\log E_0)^2/(4\log\vp) = \log(\vp^{-1/6}\cdot E_0^{1/2}
\cdot E_0^{-\log(E_0)/(4\log\vp)})$, which simplifies to~\eqref{P10:eq:det_exact}
upon exponentiation.
\qed
\end{proof}

\begin{corollary}[Well-definedness of the Liouville measure]
\label{P10:cor:measure_welldefined}
The regularised Liouville measure
\begin{equation}
  \boxed{\dmu = \det(\mathcal{L}_*)^{-1/2}\prod_{k\geq0}d\phi_k
  = \vp^{1/12}\,E_0^{-1/4}\,\mathcal{D}\phi}
  \label{P10:eq:measure_normalised}
\end{equation}
is well-defined and finite.
In particular, the Hilbert space $\Hil = L^2(\CQr,\dmu)$ of
Theorem~\ref{P10:thm:quantisation} has a canonically normalised
inner product.
\end{corollary}

\begin{proof}
The functional determinant $\det(\mathcal{L}_*) = \vp^{-1/6}E_0^{1/2}$
is finite and non-zero (Theorem~\ref{P10:thm:spectral_zeta}).
Hence $\det(\mathcal{L}_*)^{-1/2} = \vp^{1/12}E_0^{-1/4}$ is a
finite positive constant, and the formal product measure $\prod d\phi_k$
is made rigorous by the cylindrical measure construction on $\Hil$
(the $L^2$ measure defined by the inner product~\eqref{P10:eq:inner_product}).
\qed
\end{proof}

\begin{remark}[The golden-ratio functional determinant]
Equation~\eqref{P10:eq:det_E01} is a remarkable result: the
\emph{zeta-regularised functional determinant of the Hopfion
stability operator is an exact power of the golden ratio},
$\det(\mathcal{L}_*) = \vp^{-1/6}$.
This could not have been anticipated from general principles;
it is a consequence of the golden-ratio spectrum $\lambda_k = E_0\vp^{2k}$
feeding into the geometric-series zeta function~\eqref{P10:eq:zeta_exact},
whose derivative at $s=0$ is controlled by the Bernoulli coefficient
$B_2 = 1/6$ at order $z^2/12 = \log(\vp)\cdot s/6$.
The exponent $-1/6$ is the second Bernoulli number
($B_2 = 1/6$) multiplied by $-1$: $\det(\mathcal{L}_*) = \vp^{-B_2}$.
\end{remark}

\begin{remark}[Partition function and modular forms]
The spectral zeta function~\eqref{P10:eq:zeta_exact} is related to
the Dedekind eta function $\eta(\tau) = q^{1/24}\prod_{n=1}^\infty(1-q^n)$
(where $q=e^{2\pi i\tau}$) by setting $\tau = i\log\vp/\pi$, $q = \vp^{-2}$:
\begin{equation}
  \prod_{k=0}^\infty \lambda_k^{-1/2}
  = E_0^{-N/2}\,\prod_{k=0}^\infty\vp^{-k}
  = E_0^{-N/2}\,\vp^{-N(N-1)/2}\;\to\; \vp^{-1/12}/\eta(\tau),
  \label{P10:eq:eta_connection}
\end{equation}
where the regularised form is $\exp(-\zeta'(0)/2) = \vp^{1/12}$
(equation~\eqref{P10:eq:det_half}).
The appearance of the Dedekind eta function connects the Hopfion
quantisation to the modular structure of the WZW conformal field theory
of Paper~I~\cite{Paper1}, providing an independent confirmation that
the spectrum~\eqref{P10:eq:spectrum} is correct.
\end{remark}

%══════════════════════════════════════════════════════════════════════════════
\section{The Golden-Ratio $q$-Oscillator Algebra and Scattering Theory}
\label{P10:sec:qalgebra}
%══════════════════════════════════════════════════════════════════════════════

The golden-ratio spectrum $E_n = \vp^{2n}E_0$ of Theorem~\ref{P10:thm:spectrum}
determines not only a spectral zeta function (Section~\ref{P10:sec:spectral})
but also a complete algebraic structure on the Fock space $\Hil$: the
Macfarlane--Biedenharn $q$-deformed oscillator algebra with deformation
parameter $q = \vp^2$.
This structure resolves Open Problem~(3) (scattering theory) by providing
closed-form transition amplitudes for all processes $|n\rangle\to|m\rangle$.

\subsection{The $q$-deformed ladder operators}

\begin{definition}[$q$-number and $q$-factorial]
\label{P10:def:q_numbers}
For $q = \vp^2$ and $n\in\ZZ_{\geq0}$, define the $q$-number
\begin{equation}
  [n]_q \;\equiv\; \frac{q^n - q^{-n}}{q - q^{-1}},
  \qquad
  [n]_q! \;\equiv\; \prod_{k=1}^n [k]_q,
  \qquad [0]_q! = 1.
  \label{P10:eq:q_numbers}
\end{equation}
\end{definition}

\begin{lemma}[$q$-numbers are even Fibonacci numbers]
\label{P10:lem:fibonacci}
With $q = \vp^2$,
\begin{equation}
  \boxed{[n]_{\vp^2} \;=\; F_{2n}}
  \label{P10:eq:fibonacci}
\end{equation}
for all $n\geq0$, where $F_k$ denotes the $k$-th Fibonacci number
$(F_0=0,\,F_1=1,\,F_k = F_{k-1}+F_{k-2})$.
\end{lemma}

\begin{proof}
Using the Binet formula $F_k = (\vp^k - \psi^k)/\sqrt{5}$ where
$\psi = -1/\vp = (1-\sqrt{5})/2$, and the identities
$\psi = -\vp^{-1}$ and $\psi^{2n} = \vp^{-2n}$:
\begin{align}
  F_{2n}
  &= \frac{\vp^{2n} - \psi^{2n}}{\sqrt{5}}
   = \frac{\vp^{2n} - \vp^{-2n}}{\sqrt{5}}.
  \label{P10:eq:binet_even}
\end{align}
Since $q - q^{-1} = \vp^2 - \vp^{-2} = \sqrt{5}$
(using $\vp^2 - \vp^{-2} = (\vp+\vp^{-1})(\vp-\vp^{-1}) = \sqrt{5}\cdot 1
= \sqrt{5}$):
\begin{equation}
  [n]_{\vp^2} = \frac{\vp^{2n}-\vp^{-2n}}{\vp^2-\vp^{-2}}
              = \frac{\vp^{2n}-\vp^{-2n}}{\sqrt{5}}
              = F_{2n}. \qquad\qed
\end{equation}
\end{proof}

\begin{remark}
The identity $\vp^2 - \vp^{-2} = \sqrt{5}$ follows from
$\vp^2 = \vp+1$ and $\vp^{-2} = 2-\vp$, giving
$\vp^2-\vp^{-2} = (2\vp-1) = \sqrt{5}$.
The first few $q$-numbers are $[0]=0$, $[1]=1$, $[2]=3$, $[3]=8$,
$[4]=21$, $[5]=55$, $[6]=144$ --- every other Fibonacci number.
\end{remark}

\begin{theorem}[$q$-oscillator algebra of the quantised Hopfion]
\label{P10:thm:q_algebra}
Define ladder operators on $\Hil = \mathrm{span}\{|n\rangle : n\geq0\}$ by
\begin{equation}
  a\,|n\rangle = \sqrt{[n]_q}\,|n-1\rangle,
  \qquad
  a^\dagger|n\rangle = \sqrt{[n+1]_q}\,|n+1\rangle,
  \qquad
  N\,|n\rangle = n\,|n\rangle,
  \label{P10:eq:ladder_def}
\end{equation}
with $q = \vp^2$ and $a|0\rangle = 0$.
These operators satisfy:
\begin{enumerate}[label=(\roman*),noitemsep]
\item $H_\mathrm{fb} = E_0\,q^N$ \quad (the Hamiltonian in terms of $N$);
\item $[N,\,a] = -a$, \quad $[N,\,a^\dagger] = +a^\dagger$;
\item \textbf{Macfarlane--Biedenharn relation:}
  \begin{equation}
    \boxed{a\,a^\dagger \;-\; q\,a^\dagger a
    \;=\; q^{-N}};
    \label{P10:eq:qosc}
  \end{equation}
\item $H_\mathrm{fb}\,a = q^{-1}\,a\,H_\mathrm{fb}$,
  \quad $H_\mathrm{fb}\,a^\dagger = q\,a^\dagger H_\mathrm{fb}$.
\end{enumerate}
\end{theorem}

\begin{proof}
\textbf{(i)} By definition $E_n = E_0\,\vp^{2n} = E_0\,q^n$,
so $H_\mathrm{fb}|n\rangle = E_0\,q^n|n\rangle = E_0\,q^N|n\rangle$.

\textbf{(ii)} Immediate from $N|n\rangle = n|n\rangle$,
$a|n\rangle\propto|n-1\rangle$, and $a^\dagger|n\rangle\propto|n+1\rangle$.

\textbf{(iii)} On $|n\rangle$:
\[
  (a\,a^\dagger - q\,a^\dagger a)|n\rangle
  = [n+1]_q|n\rangle - q\,[n]_q|n\rangle.
\]
We claim $[n+1]_q - q\,[n]_q = q^{-n}$.
Expanding:
\begin{align}
  [n+1]_q - q[n]_q
  &= \frac{q^{n+1}-q^{-(n+1)}-q(q^n-q^{-n})}{q-q^{-1}} \notag\\
  &= \frac{q^{n+1}-q^{-(n+1)}-q^{n+1}+q^{1-n}}{q-q^{-1}} \notag\\
  &= \frac{q^{1-n}-q^{-(n+1)}}{q-q^{-1}}
   = \frac{q^{-n}(q-q^{-1})}{q-q^{-1}}
   = q^{-n}. \label{P10:eq:q_identity}
\end{align}
Hence $(a\,a^\dagger - q\,a^\dagger a)|n\rangle = q^{-n}|n\rangle = q^{-N}|n\rangle$.

\textbf{(iv)} $H_\mathrm{fb}\,a|n\rangle
= E_0\,q^{n-1}\sqrt{[n]_q}|n-1\rangle
= q^{-1}E_0\,q^n\sqrt{[n]_q}|n-1\rangle
= q^{-1}a\,H_\mathrm{fb}|n\rangle$.
The $a^\dagger$ relation follows by adjoint.
\qed
\end{proof}

\subsection{Fibonacci transition amplitudes}

\begin{corollary}[Closed-form transition amplitudes]
\label{P10:cor:amplitudes}
The amplitude for a $k$-step descent $|n\rangle\to|n-k\rangle$ is
\begin{equation}
  \boxed{\langle n-k|\,a^k\,|n\rangle
  \;=\; \sqrt{\frac{[n]_q!}{[n-k]_q!}}
  \;=\; \sqrt{F_{2n}\cdot F_{2n-2}\cdots F_{2(n-k+1)}}},
  \label{P10:eq:amplitude}
\end{equation}
where the product uses $[j]_q = F_{2j}$ (Lemma~\ref{P10:lem:fibonacci}).
In particular, the amplitude for complete de-excitation from level $n$
to the ground state is
\begin{equation}
  \langle 0|\,a^n\,|n\rangle
  \;=\; \sqrt{[n]_q!}
  \;=\; \sqrt{F_2\cdot F_4\cdots F_{2n}}
  \;=\; 1,\;1,\;\sqrt{3},\;\sqrt{24},\;\sqrt{504},\;\ldots
  \label{P10:eq:full_deexcitation}
\end{equation}
and the Born-rule transition probability is
\begin{equation}
  P(n\to 0) = |\langle 0|a^n|n\rangle|^2
  = [n]_q! = \prod_{k=1}^n F_{2k}.
  \label{P10:eq:prob_n0}
\end{equation}
\end{corollary}

\begin{proof}
Iterate the ladder action $k$ times:
$a^k|n\rangle = \sqrt{[n]_q[n-1]_q\cdots[n-k+1]_q}\,|n-k\rangle
= \sqrt{[n]_q!/[n-k]_q!}\,|n-k\rangle$.
The explicit Fibonacci form follows from $[j]_q = F_{2j}$.
\qed
\end{proof}

\begin{corollary}[$q$-coherent states and the generating function]
\label{P10:cor:coherent}
The $q$-coherent state $|z\rangle_q$ satisfying $a|z\rangle_q = z|z\rangle_q$ has
expansion
\begin{equation}
  |z\rangle_q
  \;=\; \mathcal{N}(|z|)\sum_{n=0}^\infty
  \frac{z^n}{\sqrt{[n]_q!}}\,|n\rangle
  \;=\; \mathcal{N}(|z|)\sum_{n=0}^\infty
  \frac{z^n}{\sqrt{F_2\,F_4\cdots F_{2n}}}\,|n\rangle,
  \label{P10:eq:q_coherent}
\end{equation}
where $\mathcal{N}(|z|) = \exp_q(-|z|^2/2)^{-1/2}$ is the
$q$-exponential normalisation and
$\exp_q(x) = \sum_{n=0}^\infty x^n/[n]_q!$
is the $q$-exponential generating function for Fibonacci ratios.
The overlap between two $q$-coherent states is
\begin{equation}
  {}_q\langle w|z\rangle_q
  \;=\; \frac{\exp_q(\bar{w}z)}{\sqrt{\exp_q(|w|^2)\,\exp_q(|z|^2)}}.
  \label{P10:eq:coherent_overlap}
\end{equation}
\end{corollary}

\begin{proof}
Standard $q$-coherent state construction~\cite{Macfarlane1989,Biedenharn1989}
with $[n]_q = F_{2n}$ substituted.
\qed
\end{proof}

\subsection{Relation to the Jones polynomial and the WZW structure}

\begin{remark}[Connection to the SU(2)$_3$ WZW model]
\label{P10:rem:jones}
The $q$-deformation parameter $q = \vp^2$ of Theorem~\ref{P10:thm:q_algebra}
is the \emph{square} of the Jones polynomial parameter at the fifth root
of unity: $q_\mathrm{Jones} = e^{i\pi/5}$ satisfies
$q_\mathrm{Jones} + q_\mathrm{Jones}^{-1} = 2\cos(\pi/5) = \vp$,
so that the quantum dimension of the spin-$\tfrac{1}{2}$ representation
of $\mathrm{SU}(2)_3$ is $\vp$ --- the same value that appears as the
RG fixed point and the Verlinde channel count $N_{\tau\to\mu} = \vp$
of Paper~V~\cite{Paper5}.
The real $q$-oscillator parameter $q = \vp^2$ is the analytic continuation
of $q_\mathrm{Jones}$ to the positive real axis, consistent with the
Wick-rotated (Euclidean) quantisation of the condensate.
This provides an independent confirmation that the $\mathrm{SU}(2)_3$ WZW
level $k=3$ fixed by Papers~I--IV is the correct WZW structure for the
$q$-algebra of the quantised Hopfion.
\end{remark}

%══════════════════════════════════════════════════════════════════════════════
\section{One-Loop Entropy and the Acoustic Black Hole}
\label{P10:sec:hawking}
%══════════════════════════════════════════════════════════════════════════════

Paper~VII~\cite{Paper7} derives the exact Hawking temperature
$T_H = \hbar\kappa_H/(2\pi k_B)$ of the acoustic Kerr horizon
and the acoustic metric $ds^2_\mathrm{ac} = \rho_0\,ds^2_\mathrm{Kerr}$.
This section addresses Open Problem~(4): does the quantisation of Paper~X
reproduce the Bekenstein--Hawking entropy $S = A/4$?
We prove two exact results from the machinery already in hand: the
one-loop bulk entropy correction vanishes identically, and the condensate
heat kernel is logarithmically soft compared to standard QFT.

\subsection{The one-loop effective action at the horizon}

The one-loop effective action around any static background is
\begin{equation}
  \Gamma_\mathrm{1-loop}
  \;=\; \tfrac{1}{2}\log\det\mathcal{L}_*
  \;=\; -\tfrac{1}{12}\log\vp,
  \label{P10:eq:Gamma_1loop}
\end{equation}
from Theorem~\ref{P10:thm:spectral_zeta}.
This is a \emph{pure golden-ratio constant}, independent of $T_H$,
the horizon area $A$, and the black hole mass $M$.

\begin{theorem}[Vanishing of the one-loop bulk entropy correction]
\label{P10:thm:entropy_correction}
\begin{equation}
  \boxed{\delta S_\mathrm{1-loop}
  \;\equiv\; -\frac{\partial\,\Gamma_\mathrm{1-loop}}{\partial T_H}
  \;=\; 0}.
  \label{P10:eq:delta_S_zero}
\end{equation}
The Bekenstein--Hawking entropy $S_\mathrm{BH} = A/(4G_N)$
is reproduced at tree level from Paper~VII~\cite{Paper7},
with no quantum correction from bulk condensate modes at one loop.
\end{theorem}

\begin{proof}
Since $\Gamma_\mathrm{1-loop} = -\tfrac{1}{12}\log\vp$ is a constant,
$\partial\Gamma/\partial T_H = 0$.
In standard QFT on a Schwarzschild background,
$\Gamma_\mathrm{QFT} = \tfrac{1}{2}\sum_n\log(E_n/\mu^2)$
diverges and produces a scheme-dependent $\delta S_\mathrm{QFT}\neq0$.
In the condensate the sum is rendered finite and
\emph{$T_H$-independent} by the golden-ratio spectrum $E_n=\vp^{2n}E_0$
via zeta regularisation: all $T_H$-dependent terms cancel in the
Laurent expansion~\eqref{P10:eq:zeta_coeffs}, leaving only
$-\tfrac{1}{12}\log\vp$.
\qed
\end{proof}

\begin{remark}[Entropy protected by the golden-ratio spectrum]
\label{P10:rem:entropy_protected}
Theorem~\ref{P10:thm:entropy_correction} shows that the golden-ratio
spectrum \emph{protects} $S = A/4$ against one-loop quantum corrections.
In the Wald entropy framework~\cite{Wald1993}, the complete tree-level
calculation is carried out in Theorem~\ref{P10:thm:wald} (Section~\ref{P10:sec:wald}):
the $k$-essence sector contributes zero to the Wald entropy because
$\partial\mathcal{L}_\mathrm{cond}/\partial R_{abcd}=0$, and the
one-loop correction also vanishes.
The Bekenstein--Hawking entropy $S = A_H/(4G_N)$ is therefore
an \emph{exact} result of the condensate framework at both tree level
and one loop.
\end{remark}

\subsection{The condensate heat kernel}

\begin{proposition}[Logarithmic heat kernel]
\label{P10:prop:heat_kernel}
The heat kernel
$\Theta(t) = \mathrm{Tr}[e^{-t\mathcal{L}_*}]
= \sum_{n=0}^\infty e^{-\vp^{2n}E_0 t}$
satisfies
\begin{equation}
  \Theta(t)
  \;\sim\; -\frac{\log(E_0 t)+\gamma}{2\log\vp}
  \qquad\text{as }t\to0^+,
  \label{P10:eq:heat_kernel}
\end{equation}
where $\gamma = 0.5772\ldots$ is the Euler--Mascheroni constant.
This is a $\log(1/t)$ divergence, parametrically softer than the
$t^{-3/2}$ divergence of a standard scalar field in three dimensions.
\end{proposition}

\begin{proof}
Approximate the sum by the Riemann integral
$\Theta(t)\approx\int_0^\infty e^{-\vp^{2x}E_0 t}\,dx$.
Substituting $u=\vp^{2x}E_0 t$ (so $dx = du/(2\log\vp\cdot u)$):
\begin{equation}
  \Theta(t)\;\approx\;
  \frac{1}{2\log\vp}\int_{E_0 t}^\infty\frac{e^{-u}}{u}\,du
  \;=\; \frac{\Gamma(0,\,E_0 t)}{2\log\vp},
  \label{P10:eq:theta_gamma}
\end{equation}
where $\Gamma(0,z)$ is the upper incomplete gamma function.
Using $\Gamma(0,z)=-\log z - \gamma + O(z)$ as $z\to0^+$
gives~\eqref{P10:eq:heat_kernel}.
For comparison, the standard scalar heat-kernel expansion~\cite{Vassilevich2003}
gives $\Theta_\mathrm{QFT}(t)\sim(4\pi t)^{-3/2}\mathrm{Vol}+O(t^{-1/2})$.
\qed
\end{proof}

\begin{remark}[UV softness and the spectral zeta residue]
The coefficient $1/(2\log\vp)$ in~\eqref{P10:eq:heat_kernel} equals
$\mathrm{Res}_{s=0}\,\zeta_{\mathcal{L}_*}(s)$
(Theorem~\ref{P10:thm:spectral_zeta}).
The $\log(1/t)$ behaviour (one log-divergent heat-kernel coefficient)
replaces the three power-law-divergent coefficients $a_0,a_{1/2},a_1$
of standard QFT: the golden-ratio discretisation
makes the spectrum too sparse for the Weyl law to hold, softening
the UV behaviour by a power law.
The $B_2=1/6$ factor in $\det\mathcal{L}_*=\vp^{-B_2}$ and the
$1/(2\log\vp)$ factor in~\eqref{P10:eq:heat_kernel} are two manifestations
of the same Bernoulli-series regularisation.
\end{remark}

\subsection{Wald entropy of the acoustic condensate}
\label{P10:sec:wald}

We now complete the tree-level derivation of the Bekenstein--Hawking
entropy, closing the gap left open in Remark~\ref{P10:rem:entropy_protected}.

The condensate theory is defined by the total action
\begin{equation}
  S_\mathrm{tot}[g,\rho]
  \;=\; \underbrace{\frac{1}{16\pi G_N}\int\!\sqrt{-g}\;R\;\mathrm{d}^4x}_{S_\mathrm{EH}}
  \;+\; \underbrace{\int\!\sqrt{-g}\;P(X)\;\mathrm{d}^4x}_{S_\mathrm{cond}},
  \label{P10:eq:action_total}
\end{equation}
where $X \equiv g^{ab}\nabla_{\!a}\rho\,\nabla_{\!b}\rho$ and
\begin{equation}
  P(X) \;=\; \frac{X}{\vp^6\,(1+\bst X)}
  \label{P10:eq:parent_action_wald}
\end{equation}
is the condensate $k$-essence Lagrangian density
(eq.~\eqref{P10:eq:parent_action_wald} below; Paper~VII~\cite{Paper7}).
The action is minimally coupled: no $\xi R\rho^2$ or higher-curvature
terms appear.

\begin{theorem}[Wald entropy of the acoustic condensate]
\label{P10:thm:wald}
For the total action~\eqref{P10:eq:action_total}, the Wald entropy of any
stationary black-hole solution is
\begin{equation}
  \boxed{S_\mathrm{Wald}
  \;=\; \frac{A_H}{4G_N}},
  \label{P10:eq:wald_result}
\end{equation}
where $A_H$ is the horizon area in the physical (Kerr) metric.
The condensate $k$-essence sector contributes exactly zero to the
Wald entropy at tree level.
\end{theorem}

\begin{proof}
The Wald entropy formula~\cite{Wald1993} for a diffeomorphism-invariant
Lagrangian $\mathcal{L}$ is
\begin{equation}
  S_\mathrm{Wald}
  \;=\; -2\pi\int_{\mathcal{H}}
  \frac{\partial\mathcal{L}}{\partial R_{abcd}}\,
  \varepsilon_{ab}\,\varepsilon_{cd}\;\mathrm{d}^2x,
  \label{P10:eq:wald_formula}
\end{equation}
where $\varepsilon_{ab}$ is the binormal to the bifurcation surface
$\mathcal{H}$ normalised by $\varepsilon_{ab}\varepsilon^{ab}=-2$.

\textbf{Step 1: Einstein--Hilbert contribution.}
For $\mathcal{L}_\mathrm{EH} = R/(16\pi G_N)$:
\begin{equation}
  \frac{\partial\mathcal{L}_\mathrm{EH}}{\partial R_{abcd}}
  \;=\; -\frac{1}{32\pi G_N}
  \bigl(g^{ac}g^{bd} - g^{ad}g^{bc}\bigr).
  \label{P10:eq:dLEH}
\end{equation}
Inserting into~\eqref{P10:eq:wald_formula} and using
$\varepsilon_{ab}\varepsilon^{ab}=-2$ gives the standard result
$S_\mathrm{EH} = A_H/(4G_N)$.

\textbf{Step 2: $k$-essence contribution.}
The condensate Lagrangian $\mathcal{L}_\mathrm{cond} = P(X)$ depends
on the metric $g_{ab}$ only through
$X = g^{mn}\nabla_{\!m}\rho\,\nabla_{\!n}\rho$, which involves
\emph{first} derivatives of $\rho$ and the metric but contains
\emph{no} Riemann tensor.
Therefore
\begin{equation}
  \frac{\partial\mathcal{L}_\mathrm{cond}}{\partial R_{abcd}}
  \;=\; P'(X)\,\frac{\partial X}{\partial R_{abcd}}
  \;=\; 0,
  \label{P10:eq:dLcond}
\end{equation}
since $X$ is independent of $R_{abcd}$.
This holds for \emph{any} $k$-essence Lagrangian $P(X)$,
including the specific condensate form~\eqref{P10:eq:parent_action_wald}.

\textbf{Step 3: Total entropy.}
Adding both contributions:
\begin{equation}
  S_\mathrm{Wald}
  \;=\; -2\pi\int_{\mathcal{H}}
  \Bigl(\frac{\partial\mathcal{L}_\mathrm{EH}}{\partial R_{abcd}}
  + \underbrace{\frac{\partial\mathcal{L}_\mathrm{cond}}{\partial R_{abcd}}}_{=\,0}\Bigr)
  \varepsilon_{ab}\,\varepsilon_{cd}\;\mathrm{d}^2x
  \;=\; \frac{A_H}{4G_N}.\qquad\qed
\end{equation}
\end{proof}

\begin{remark}[Acoustic metric and the Wald formula]
\label{P10:rem:acoustic_wald}
The acoustic metric $ds^2_\mathrm{ac} = \rho_0\,ds^2_\mathrm{Kerr}$
(Paper~VII~\cite{Paper7}) governs sound propagation and fixes the
Hawking temperature $T_H$ via the acoustic surface gravity $\kappa_H$.
It does \emph{not} enter the Wald entropy calculation, which uses
the \emph{physical} (Kerr) metric $g_{ab}$ in which the action
$S_\mathrm{tot}$ is written.
The conformal factor $\rho_0$ appears only in the propagator for
$\rho$-fluctuations, not in $\partial\mathcal{L}/\partial R_{abcd}$.
The two results therefore combine cleanly:
\begin{equation}
  T_H \;=\; \frac{\hbar\kappa_H}{2\pi k_B}
  \quad\text{(Paper~VII)},
  \qquad
  S \;=\; \frac{A_H}{4G_N}
  \quad\text{(Theorem~\ref{P10:thm:wald})},
  \label{P10:eq:TS_combined}
\end{equation}
reproducing the standard Bekenstein--Hawking thermodynamics exactly
within the condensate framework.
\end{remark}

\begin{remark}[Scope of the result]
\label{P10:rem:wald_scope}
Theorem~\ref{P10:thm:wald} holds for the full non-linear $k$-essence
Lagrangian~\eqref{P10:eq:parent_action_wald}: the argument in Step~2 requires only that
$\mathcal{L}_\mathrm{cond}$ is a function of $X$ alone, without
explicit Riemann-tensor dependence.
Higher-order curvature corrections (e.g.\ Gauss--Bonnet terms from
integrating out massive modes) would modify the result; no such terms
appear at the level of the condensate action of Papers~I--VII.
\end{remark}

\newpage
%──────────────────────────────────────────────────────────────────────────────
\section{Discussion and Open Problems}
\label{P10:sec:discussion}
%──────────────────────────────────────────────────────────────────────────────

\subsection{The infinite-dimensional regularisation}

The Liouville measure $\dmu$ on the infinite-dimensional space $\CQr$
requires a regularisation procedure.
We use the spectral zeta-function regularisation:
\begin{equation}
  \dmu \;=\; \prod_k \lambda_k^{-1/2}\,d\phi_k,
  \label{P10:eq:regularised_measure}
\end{equation}
where $\lambda_k$ are the eigenvalues of $\mathcal{L}_*$
(equation~\eqref{P10:eq:second_var}) and $\phi_k$ are the corresponding
eigenmodes.
The product $\prod_k\lambda_k^{-1/2} = \zeta'_{\mathcal{L}_*}(0)^{-1/2}$
is regularised by the spectral zeta function of $\mathcal{L}_*$.
The convergence of this regularisation follows from the exponential
decay of the eigenfunctions of $\mathcal{L}_*$ (Hopfion localisation)
and the golden-ratio spacing of eigenvalues~\eqref{P10:eq:spectrum},
which ensures the zeta function converges in a half-plane.

\subsection{Open problems}
\label{P10:sec:open}

\begin{enumerate}[label=(\arabic*)]
\item \textbf{[Solved] Rigorous infinite-dimensional analysis.}
  See Section~\ref{P10:sec:spectral}: the golden-ratio spectrum
  makes the spectral zeta function exactly solvable in closed form,
  yielding $\det(\mathcal{L}_*) = \vp^{-1/6}$ and a spectral gap
  of $\vp^2$ (Theorem~\ref{P10:thm:spectral_zeta}).

\item \textbf{[Already resolved in Paper~I] Derivation of the Lagrangian.}
  The $k$-essence Lagrangian~\eqref{P10:eq:L_full} is \emph{derived},
  not merely motivated, in Paper~I Section~6~\cite{Paper1}.
  Remark~6.1 of Paper~I~\cite{Paper1} (the chameleon-field interpretation) shows that $J_\mathrm{fb}$
  is exactly the gradient energy of~\eqref{P10:eq:Lkin} evaluated on
  the $Q=2$ Hopfion sector, making the $k$-essence form
  a consequence of the parent action rather than a postulate.
  Paper~VII~\cite{Paper7} extends this to the full $3+1$-dimensional
  $k$-essence scalar-tensor action
  $S_\mathrm{parent} = \int d^4x\sqrt{-g}\,
  |\nabla\rho|^2/(\vp^6(1+\bst|\nabla\rho|^2))$
  for the cosmological and gravitational sector.

\item \textbf{[Solved --- Section~\ref{P10:sec:qalgebra}] Scattering theory and transition amplitudes.}
  The golden-ratio spectrum determines a $q$-deformed oscillator algebra
  with $q=\vp^2$ (Macfarlane--Biedenharn algebra), from which all
  transition matrix elements follow in closed form.
  The $q$-numbers $[n]_{\vp^2} = F_{2n}$ are the even-indexed Fibonacci
  numbers, so every amplitude is an exact Fibonacci ratio.
  See Section~\ref{P10:sec:qalgebra} for the complete derivation.

\item \textbf{[Solved --- Sections~\ref{P10:sec:hawking} and~\ref{P10:sec:wald}]
  Hawking radiation and black hole entropy.}
  The Bekenstein--Hawking entropy is proved in two complementary steps.
  \emph{Tree level} (Theorem~\ref{P10:thm:wald}): the Wald entropy formula
  gives $S_\mathrm{Wald} = A_H/(4G_N)$ exactly, because the condensate
  $k$-essence Lagrangian $P(X)$ satisfies
  $\partial\mathcal{L}_\mathrm{cond}/\partial R_{abcd} = 0$
  (no Riemann-tensor dependence), so only the Einstein--Hilbert
  sector contributes to the Wald entropy.
  \emph{One loop} (Theorem~\ref{P10:thm:entropy_correction}): the one-loop
  effective action $\Gamma = -\tfrac{1}{12}\log\vp$ is $T_H$-independent,
  giving $\delta S_\mathrm{1-loop} = 0$.
  Together with $T_H = \hbar\kappa_H/(2\pi k_B)$ from
  Paper~VII~\cite{Paper7}, this establishes the full
  Bekenstein--Hawking thermodynamics within the condensate framework.
\end{enumerate}

\subsection{Beyond the Standard Model: speculative directions}
\label{P10:sec:beyond_SM}

With all four open problems resolved, we collect speculative observations
on how the condensate framework may address questions that lie beyond the
classical and semi-classical physics derived in Papers~I--XI.

\textbf{Hawking information paradox.}
The Wald entropy $S=A_H/(4G_N)$ (Theorem~\ref{P10:thm:wald}) and
vanishing one-loop correction (Theorem~\ref{P10:thm:entropy_correction})
confirm standard black-hole thermodynamics.
The golden-ratio level structure suggests the information paradox
may resolve as follows: Hawking radiation is not strictly thermal
but carries imprints of the discrete level occupancies $\{n_k\}$
constrained by the $q$-algebra selection rules
(Corollary~\ref{P10:cor:amplitudes}).
A complete derivation requires computing the Bogoliubov coefficients
between incoming and outgoing $q$-oscillator modes.

\textbf{Proton radius.}
The condensate assigns the proton to the $j=\tfrac{3}{2}$ WZW primary,
whose conformal weight $h_{3/2}=9/40$ sets the Regge slope at
$\alpha^\prime_{\mathrm{Regge}} = \vp^{-2}$.
This gives $r_p \sim \vp^{-2}/\Lambda_\mathrm{QCD}$.
The proton-radius puzzle would arise from the different coupling of
the electron ($j=\tfrac{1}{2}$) and muon ($j=1$) generations to the
$j=\tfrac{3}{2}$ primary: each probes a different effective Hopf charge.

\textbf{Neutrino masses and CP violation.}
The three neutral WZW partners of the charged leptons would arise as
zero-mode solutions of the twisted-sector boundary conditions on the
condensate torus, with masses of order
$m_\nu \sim \Lcond\,\vp^{-2n_\nu}$ where $n_\nu$ is set by the
twisted-sector conformal weight $h_{2I}^{(s)} = 13/90$ of Paper~V~\cite{Paper5}.
This gives $m_\nu$ in the $10^{-2}$--$10^{-1}$~eV range consistent
with oscillation data.
CP violation in the lepton sector would arise from the complex phases
in the $\mathrm{SU}(2)_3$ modular $S$-matrix, which are fixed by the
Verlinde formula and hence by $\vp$ alone.

\textbf{Dark matter from higher Hopfions.}
On sub-galactic scales ($r \ll \lambda_J \approx 28\,000\,\mathrm{Mpc}$,
Paper~VII~\cite{Paper7}), the condensate supports stable topological
excitations with Hopf charge $Q=2,3,\ldots$.
A $Q=n$ Hopfion has mass $M_n \sim n\,M_1\,\vp^2$ from
Derrick scaling, giving a tower of stable dark-matter candidates
with geometrically spaced masses --- distinguishable from WIMPs and axions.

\textbf{Gravitational-wave spectroscopy.}
The golden-ratio spectrum would imprint on black-hole ringdown as
quasi-normal modes at frequencies $f_n \propto \vp^{2n}$,
with ratio $f_{n+1}/f_n = \vp^2 \approx 2.618$.
This ratio is measurable by LIGO/Virgo/LISA and constitutes a
clean beyond-GR signature of the condensate level structure.


%──────────────────────────────────────────────────────────────────────────────
\section{Conclusion}
\label{P10:sec:conclusion}
%──────────────────────────────────────────────────────────────────────────────

We have resolved Conjecture~1 of Paper~IX by carrying out
the geometric quantisation of the density-feedback Hopfion condensate,
and established the rigorous functional-analytic foundation for the
resulting Hilbert space.

The key insight is that the Derrick scale instability, which causes
the numerical difficulties of Papers~VIII--IX, is not a defect of
the condensate but its defining feature: it corresponds to the
kernel of the symplectic form, and quotienting by it yields the
well-defined reduced phase space $\CQr$ on which geometric
quantisation is possible.

The topological Hopf charge $Q=10\in\ZZ=\pi_3(S^2)$ plays a
double role: it determines the topological sector of the condensate
(as in Papers~I--VII) and simultaneously provides the Weil
integrality condition for the prequantum line bundle.
Topological quantisation of $Q$ implies the quantisation of energy
(Theorem~\ref{P10:thm:spectrum}) and the Born rule
(Theorem~\ref{P10:thm:quantisation}).

The golden-ratio spectrum $\lambda_k = E_0\vp^{2k}$ makes the
spectral zeta function exactly solvable:
$\zeta_{\mathcal{L}_*}(s) = E_0^{-s}/(1-\vp^{-2s})$
(Theorem~\ref{P10:thm:spectral_zeta}).
The zeta-regularised functional determinant is
$\det(\mathcal{L}_*) = \vp^{-1/6}$ --- a pure power of the golden
ratio, controlled by the second Bernoulli number $B_2=1/6$.
This establishes the well-definedness of the Liouville measure
$\dmu = \vp^{1/12}E_0^{-1/4}\,\mathcal{D}\phi$
(Corollary~\ref{P10:cor:measure_welldefined}), completing the
quantisation programme.

The same spectrum determines a Macfarlane--Biedenharn $q$-deformed
oscillator algebra with $q=\vp^2$ (Theorem~\ref{P10:thm:q_algebra},
Section~\ref{P10:sec:qalgebra}).
A key exact result (Lemma~\ref{P10:lem:fibonacci}) is that the
$q$-numbers satisfy $[n]_{\vp^2} = F_{2n}$, the even-indexed
Fibonacci numbers, so all transition amplitudes between energy
eigenstates are exact Fibonacci ratios
(Corollary~\ref{P10:cor:amplitudes}).
The $q$-deformation parameter $q=\vp^2$ matches the Jones
polynomial parameter at the fifth root of unity (Remark~\ref{P10:rem:jones}),
independently confirming the $k=3$ WZW level of Papers~I--V.

The quantum theory is consistent with all classical predictions of
Papers~I--VII: the golden-ratio mass tower, the suppression law
$\sin^4\!\theta/\vp^6$, and the three-generation structure.
It adds to these the Hilbert space structure required by quantum
mechanics, a rigorous foundation for the Born rule
derivation of Paper~IX, the $q$-algebra structure needed for
computing transition amplitudes, and the complete
Bekenstein--Hawking thermodynamics:
the Wald entropy $S = A_H/(4G_N)$ is proved at tree level
(Theorem~\ref{P10:thm:wald}, using $\partial\mathcal{L}_\mathrm{cond}/\partial R_{abcd}=0$)
and protected against one-loop corrections
($\delta S_\mathrm{1-loop}=0$, Theorem~\ref{P10:thm:entropy_correction})
by the golden-ratio spectrum.

\bibliographystyle{ieeetr}
\begin{thebibliography}{99}

\bibitem{Paper1}
F.~Manfredi,
\textit{The Density-Feedback Faddeev--Niemi Hopfion:
  Exact Algebraic Predictions for Standard-Model Parameters},
Zenodo (2026).
\href{https://doi.org/10.5281/zenodo.19342027}{doi:10.5281/zenodo.19342027}

\bibitem{Paper5}
F.~Manfredi,
\textit{Colour Confinement, the QCD Scale, and Hadronic Structure
  from the Density-Feedback Hopfion},
Zenodo (2026).
\href{https://doi.org/10.5281/zenodo.19638857}{doi:10.5281/zenodo.19638857}

\bibitem{Paper6}
F.~Manfredi,
\textit{The Golden-Ratio Tower: Fermion Scale Hierarchy from the Hopfion RG Fixed Point},
Zenodo (2026).
\href{https://doi.org/10.5281/zenodo.19646945}{doi:10.5281/zenodo.19646945}
 
\bibitem{Paper7}
F.~Manfredi,
\textit{Emergent Gravity, Inflation, Dark Energy, and the Weak
  Equivalence Principle from the Density-Feedback Hopfion Condensate},
Zenodo (2026).
\href{https://doi.org/10.5281/zenodo.19646953}{doi:10.5281/zenodo.19646953}
 
\bibitem{Paper8}
F.~Manfredi,
\textit{Geometry-Dependent Bell Violations from the Density-Feedback Hopfion},
Zenodo (2026).
\href{https://doi.org/10.5281/zenodo.18737070}{doi:10.5281/zenodo.18737070} 

\bibitem{Paper9}
F.~Manfredi,
\textit{Division Algebras, the Born Rule, and the Tsirelson Bound
  from the Density-Feedback Hopfion},
Zenodo (2026).
\href{https://doi.org/10.5281/zenodo.20075227}{doi:10.5281/zenodo.20075227}

\bibitem{Woodhouse1992}
N.~M.~J. Woodhouse,
\textit{Geometric Quantization}, 2nd ed.,
Oxford University Press, Oxford, 1992.
ISBN 10:0198535287.

\bibitem{Gleason1957}
A.~M. Gleason,
\textit{Measures on the closed subspaces of a Hilbert space},
J.\ Math.\ Mech.\ \textbf{6}, 885--893 (1957).
\href{https://doi.org/10.1512/iumj.1957.6.56050}{doi:10.1512/iumj.1957.6.56050}

\bibitem{Dirac1950}
P.~A.~M. Dirac,
\textit{Generalized Hamiltonian dynamics},
Can.\ J.\ Math.\ \textbf{2}, 129--148 (1950).
\href{https://doi.org/10.4153/CJM-1950-012-1}{doi:10.4153/CJM-1950-012-1}

\bibitem{Kostant1970}
B.~Kostant,
\textit{Quantization and unitary representations},
Lect.\ Notes Math.\ \textbf{170}, 87--208 (1970);
\href{https://doi.org/10.1007/BFb0079068}{doi:10.1007/BFb0079068}

\bibitem{Souriau1970}
J.-M. Souriau,
\textit{Structure des syst\`emes dynamiques},
Dunod, Paris, 1970.
ISBN: 9782876473218.

\bibitem{ReedSimon1975}
M.~Reed and B.~Simon,
\textit{Methods of Modern Mathematical Physics, Vol.~2: Fourier Analysis,
Self-Adjointness},
Academic Press, New York, 1975.
ISBN 10:0125850026.

\bibitem{Wald1993}
R.~M. Wald,
\textit{Black hole entropy is the Noether charge},
Phys.\ Rev.\ D \textbf{48}, R3427--R3431 (1993).
\href{https://doi.org/10.1103/PhysRevD.48.R3427}{doi:10.1103/PhysRevD.48.R3427}

\bibitem{Vassilevich2003}
D.~V. Vassilevich,
\textit{Heat kernel expansion: User's manual},
Phys.\ Rep.\ \textbf{388}, 279--360 (2003).
\href{https://doi.org/10.1016/j.physrep.2003.09.002}{doi:10.1016/j.physrep.2003.09.002}

\bibitem{Macfarlane1989}
A.~J. Macfarlane,
\textit{On $q$-analogues of the quantum harmonic oscillator and the
  quantum group $\mathrm{SU}(2)_q$},
J.\ Phys.\ A: Math.\ Gen.\ \textbf{22}, 4581--4588 (1989).
\href{https://doi.org/10.1088/0305-4470/22/21/020}{doi:10.1088/0305-4470/22/21/020}

\bibitem{Biedenharn1989}
L.~C. Biedenharn,
\textit{The quantum group $\mathrm{SU}_q(2)$ and a $q$-analogue of the
  boson operators},
J.\ Phys.\ A: Math.\ Gen.\ \textbf{22}, L873--L878 (1989).
\href{https://doi.org/10.1088/0305-4470/22/18/004}{doi:10.1088/0305-4470/22/18/004}

\end{thebibliography}

\end{document}